\documentclass[10pt]{article}
\usepackage[a4paper,margin=0.82in]{geometry}
\usepackage{amsmath,amssymb,booktabs,microtype}
\usepackage[colorlinks=true,linkcolor=blue,urlcolor=blue]{hyperref}
\usepackage[T1]{fontenc}
\usepackage{lmodern}

\title{A finite near-block band reproduces the spectral\newline
failure census of a prime-shell operator}
\author{Liang Wang\\
School of Mathematics and Statistics\\
Huazhong University of Science and Technology (HUST)\\
Wuhan, China}
\date{September 3, 2026}

\begin{document}
\maketitle

\begin{abstract}
We test a predeclared near-block truncation suggested by the preceding
eigenmode profile audit.  A count-2048 response-blind prime-shell operator is
split into eight contiguous blocks of length 256, and the band of block
distances at most three is retained under the same full-window square-energy
normalization as the parent matrix.  On the complete panel of three origins,
three shell scales, two beta values, and 18 rows, the band reproduces exactly
the six beta=2 spectral-cap failure keys of the full matrix.  On those six
rows the selected full-mode absolute Rayleigh retention is at least
0.9915711764, while the omitted tail contributes at most 0.0084288236 in
absolute Rayleigh fraction.  The result is a finite, independently replayed
near-block reduction; it is not a causality theorem, a bandwidth-uniform
estimate, an arithmetic gain, or a twin-prime result.
\end{abstract}

\section{Question and frozen panel}

The previous audit decomposed the selected extremal eigenmode into fixed
block-distance layers.  On the six beta=2 rows that exceeded the working
spectral cap, distances zero through three carried at least 99.157 percent
of the absolute Rayleigh mass.  The next smallest operator-level question is
whether the corresponding fixed band reproduces the parent failure census.

The interval origins are
\(1010001,1018021,1026041\).  Each interval has 2048 consecutive integers,
partitioned into eight contiguous blocks of length 256.  We use
\(Q\in\{512,2048,8192\}\), kernel exponent one, the all-plus law, and
\(\beta\in\{0,2\}\), giving 18 rows.  The complete Cartesian panel is
materialized before any spectral result is read.  The working caps are
0.64 for the spectral norm and 0.83 for the Schur row-sum envelope.

\section{Operator and exact truncation identity}

For \(Q<p\leq 2Q\), define
\[
B_p(u,t)=p\frac{66^2}{66^2+(u-t)^2}
\left({\bf1}_{p\mid u-t}-\frac{1}{p-1}\right)
{\bf1}_{u\ne t}{\bf1}_{p\nmid u}{\bf1}_{p\nmid t}.
\]
With \(w_p=(p/Q)^\beta\), let
\[
A(u,t)=\sum_p w_pB_p(u,t),\qquad
G(u)=\sum_p\sum_{s\in I}(w_pB_p(u,s))^2,
\qquad T(u,t)=\frac{A(u,t)}{\sqrt{G(u)G(t)}}.
\]
The geometry \(G\) is always computed on the full window, including for
the truncated matrix.

Writing \(b(i)=\lfloor i/256\rfloor\), freeze
\[
B_3(i,j)={\bf1}_{|b(i)-b(j)|\leq3}T(i,j),\qquad
R_3=T-B_3.
\]
The mask and its complement are disjoint and exhaustive, hence
\[
                 T=B_3+R_3                                      \tag{1}
\]
entrywise on every finite row.  Both matrices are symmetric.  If \(v\) is
the unit eigenvector of \(T\) selected by the largest absolute eigenvalue
rule (the minimum mode wins an exact tie), then
\[
 v^{\mathsf T}B_3v+v^{\mathsf T}R_3v
 =v^{\mathsf T}Tv=\lambda.                                      \tag{2}
\]
Equation (2) is an exact finite identity, not an attribution statement.

\section{Certification protocol}

The producer sums the prime shell in ascending order and writes canonical
JSON containing all 18 rows.  A separate checker has its own sieve, sums the
shell in descending order, rebuilds the full and band eigensystems, and
checks metrics, mode residuals, the Rayleigh identity, failure keys, and the
inherited rational anchor.  An adversarial checker mutates protocol,
provenance, row census, band definition, numerical fields, audit counts,
anchor, firewall, and clue fields; every mutation must be rejected.

The anchor is the exact interval \([1010346,1010359)\), at \(Q=4\), exponent
one, where the shell is \(\{5,7\}\).  It checks rational symmetry and
positive geometry and does not select a main-panel row.  Local Bridge-B
repeats the producer, independent checker, and stress suite in normal and
optimized Python modes, requiring empty standard error and byte-identical
output.  The official Route-A/Route-B evaluator files named by the Session
are absent from this checkout, so no official route pass is asserted.

\section{Results}

Table~\ref{tab:census} gives the full and band cap-failure census.  The band
preserves the six beta=2 failures exactly: all three origins at
\(Q=2048,8192\).  It also preserves the absence of beta=2 Schur failures.

\begin{table}[h]
\centering
\caption{Finite cap census on the complete 18-row panel.}
\label{tab:census}
\begin{tabular}{@{}lrrrrrr@{}}
\toprule
 & \multicolumn{2}{c}{full, beta 0} & \multicolumn{2}{c}{full, beta 2}
 & \multicolumn{2}{c}{band, beta 2}\\
quantity & spectral & Schur & spectral & Schur & spectral & Schur\\
\midrule
rows & 9 & 9 & 9 & 9 & 9 & 9\\
failures & 9 & 9 & 6 & 0 & 6 & 0\\
\bottomrule
\end{tabular}
\end{table}

For beta=2, the band spectral values at \(Q=512\) are
0.5152073962--0.5152083701, at \(Q=2048\) are
0.7037611538--0.7037762222, and at \(Q=8192\) are
0.7051585288--0.7051589866.  Thus the fixed band is below the cap at the
smallest scale and above it at the same six high-scale settings as the full
matrix.  At \(Q=512\), its value is slightly larger than the full value,
so truncation is not a monotone norm-repair operation.

\begin{table}[h]
\centering
\caption{Selected full-mode retention on the six beta=2 failure rows.}
\label{tab:retention}
\begin{tabular}{@{}lrr@{}}
\toprule
quantity & minimum & maximum\\
\midrule
absolute band Rayleigh retention & 0.9915711764 & 0.9915735754\\
absolute tail Rayleigh fraction & 0.0084264246 & 0.0084288236\\
tail Schur norm & 0.0119462342 & 0.0119686969\\
tail Frobenius norm & 0.0141458319 & 0.0141740682\\
\bottomrule
\end{tabular}
\end{table}

For comparison, over all nine beta=2 rows the band absolute retention ranges
from 0.9915711764 to 1.0016823597.  The value above one is possible because
the tail Rayleigh term can have the opposite sign to the selected negative
eigenvalue; it is another reason not to interpret the band as a monotone
approximation theorem.  The finite eigensystem residual and norm checks are
below the declared certificate tolerances, and the exact band/tail identity
holds by construction.

\section{Interpretation and claim boundary}

The main positive result is an operator-level finite reproduction: a band
chosen before the panel is evaluated has the same beta=2 spectral failure
keys as the full matrix.  Together with the preceding mode profile, this
locates the finite excess in a near-block subspace at this partition scale.
It does not show that the near-block entries cause the excess.  In
particular, a Rayleigh quotient of a full-matrix eigenvector is not an
operator-norm estimate for \(B_3\), and the small tail on this panel does not
imply a decay law for growing windows or other origins.

The next finite question is bandwidth stability: test smaller predeclared
cutoffs, retaining the same full-window normalization and caps, before any
asymptotic interpretation.

\begin{verbatim}
TPC374_BAND_REPLAY = NUMERICALLY_CERTIFIED_FINITE_18_ROWS
TPC374_BAND_FAILURE_CENSUS = NUMERICALLY_CERTIFIED_FINITE_SCOPED
TPC374_PARENT_FAILURE_REPRODUCTION = NUMERICALLY_CERTIFIED_FINITE_SCOPED
TPC374_BAND_OPERATOR_UNIFORMITY = OPEN
TPC374_CROSS_BLOCK_CAUSALITY = OPEN
TPC374_ARITHMETIC_ADVANCE = NO
TPC374_FIXED_POWER_CREDIT = 0
TPC374_FULL_GATE_B = OPEN
TPC374_TWIN_PRIME_RESULT = NONE
\end{verbatim}

No source-uniform arithmetic $L^2$ estimate, growing operator bound,
prime-shell reassembly, official Route-A/Route-B closure, or twin-prime
theorem is claimed.

\end{document}
